\section{はじめに} \label{sec1}

近年，Embodied AI~\cite{duan2022}と呼ばれる，身体性を持つエージェントが現実世界や仮想空間で学習を行い，多様な物体が存在する複雑な環境でさまざまなタスクを遂行しようとする研究が行われている．例えば，「エージェントが人間のある指示に従って家事（例：掃除，洗濯）をする」，「エージェントが環境の中にある特定の物を見つけて案内を行う」といったタスクが挙げられる．この研究分野では，自然言語の指示に従って，エージェントがタスクを実行するためのタスクプランニング（行動計画）をするために，大規模言語モデル（Large Language Models, LLMs）を活用することが注目されている~\cite{huang22,rana2023sayplan,song2023llm}．これらの研究により，LLMは日常生活でタスクを実行する上で重要な常識知識を有することが実証されてきた．

LLMを用いて，環境の状況（オブジェクト（物体）の状態や位置，物体間の位置関係）に応じてエージェントの行動を計画する研究も進められている~\cite{lin2023grounded}．
このような行動計画をLLMによって実現するためには，LLMに環境の状況を正確に認識させることが課題となる．しかし，環境が広範・複雑であるほど情報量が多くなるため，すべての情報をLLMに認識させることは非現実的である．そのため，必要な情報を絞り込むことが課題となる．

本研究では，既存の家庭内タスクベンチマーク~\cite{Puig_2018_CVPR,ALFRED20}を補完するため，同じ指示に対して初期状態に応じた行動変更が必要となる設定を診断対象とする．例えば``Turn on all lightswitches''では，同名の物体をIDで区別し，ON済みの物体とOFFの物体を追跡する必要がある．すべてONなら行動せずに終了する判断も必要となる．個々の現在状態は環境知識から取得する．

本研究では，物体の状態・ID・関係を家庭環境知識としてLLMに提示し，行動を逐次生成する実装を扱う．知識の抽出・自然言語化はコードで行い，この抽出自体にはLLMを用いない．LLMが生成候補の前提条件を充足と判定すれば実行し，棄却すれば別要求で未充足条件を尋ね，1回再生成して再検証せず実行する．ここで，行動計画はタスク遂行のための一連の行動，アクションステップは各個別の行動を指す．

提案手法の評価を行うために，本研究では，家庭内での活動を三次元仮想空間内でシミュレートするマルチエージェントプラットフォームVirtualHome~\cite{Puig_2018_CVPR}を用いて，家庭環境の把握が不可欠な家庭内タスクのデータセットを作成する．具体的には，物体の状態に着目した「状態変化タスク」と，物体間の関係に着目した「配置タスク」の2種類を対象とする．これらのタスクでは，コンピュータビション分野の研究のように，エージェント視点の視覚情報を処理して環境を認識する手法~\cite{song2023llm,kim2023context,nakajima2024combining}ではなく，シミュレータの内部から取得したデータを家庭環境知識として用いることを前提とする．また，本研究では，家庭環境知識をLLMに認識させることで，家庭環境知識に基づいた適切な行動計画を生成する手法を提案する．

本論文の貢献は以下の通りである．
\begin{itemize}
    \item 定型指示の下で同名物体の識別と初期状態に応じた行動変更を扱う，状態変化・配置の計618インスタンスの構成と，対象範囲・分割・参照手順の明示．
    \item 状態・ID・支持／包含関係の抽出と自然言語化，事例検索，LLMによる前提条件判定，1回の再生成，停止・採点を追跡できる具体的実装の仕様化．
    \item 基本性能評価，構成要素比較，ケース分析を通じて，検証方式の効果のモデル依存性と，条件判定・終了判断・実行制約に関わる失敗を具体化し，LLM逐次計画器の改善対象を処理段階別に示す診断的知見．
\end{itemize}

評価は，まず全415タスクに対する基本性能を調べ，次に4タスク・2モデル・10条件・3反復の240エピソードで構成要素の働きを比較し，実行ログとケース分析から結果を解釈する流れで行う．構成要素比較でLLM検証あり／なしの成功数はGPT-4o miniで9/12対8/12，GPT-4oで9/12対9/12だった．各分母12は目的選択した4タスクの各3反復である．

本論文の構成は次の通りである．第~\ref{sec2}では，本研究の関連研究について述べる．第~\ref{sec3}では，タスク設定とデータセット構築について述べる．第~\ref{sec4}では，本研究の提案手法とその構成について述べる．第~\ref{sec5}では，評価実験，結果，考察について述べる．最後に第~\ref{sec6}では，本研究の結論と今後の課題について述べる．

なお，本論文は，青山らの研究成果~\cite{青山2024,Aoyama2025}を発展させ，基本性能と各構成要素の評価をまとめたものである．
提案手法のソースコードとデータセットはGitHubリポジトリにまとめている~\footnote{\url{https://github.com/takanori-ugai/Precheck-LLM-Planner}}．実験条件，評価手順，結果の解釈に必要な集計値は本文に示す．コード全文，データ一覧，各試行のログと詳細集計は同リポジトリで公開する．
